%%%%%%%%%%%%%%%%%%%%%%%%%%%%%%%%%%%%%%%%%
% Short Sectioned Assignment LaTeX Template Version 1.0 (5/5/12)
% This template has been downloaded from: http://www.LaTeXTemplates.com
% Original author:  Frits Wenneker (http://www.howtotex.com)
% License: CC BY-NC-SA 3.0 (http://creativecommons.org/licenses/by-nc-sa/3.0/)
%%%%%%%%%%%%%%%%%%%%%%%%%%%%%%%%%%%%%%%%%

% \documentclass[paper=a4, fontsize=11pt]{scrartcl} % A4 paper and 11pt font size
\documentclass[11pt, a4paper]{book}
\usepackage[T1]{fontenc} % Use 8-bit encoding that has 256 glyphs
\usepackage[utf8]{inputenc}
\usepackage{fourier} % Use the Adobe Utopia font for the document - comment this line to return to the LaTeX default
\usepackage{listings} % para insertar código con formato similar al editor
\usepackage[spanish, es-tabla]{babel} % Selecciona el español para palabras introducidas automáticamente, p.ej. "septiembre" en la fecha y especifica que se use la palabra Tabla en vez de Cuadro
\usepackage{url} % ,href} %para incluir URLs e hipervínculos dentro del texto (aunque hay que instalar href)
\usepackage{graphics,graphicx, float} %para incluir imágenes y colocarlas
\usepackage[gen]{eurosym} %para incluir el símbolo del euro
\usepackage{cite} %para incluir citas del archivo <nombre>.bib
\usepackage{enumerate}
\usepackage{hyperref}
\usepackage{graphicx}
\usepackage{tabularx}
\usepackage{booktabs}
\usepackage{subfigure}

\usepackage[table,xcdraw]{xcolor}
\hypersetup{
	colorlinks=true,	% false: boxed links; true: colored links
	linkcolor=black,	% color of internal links
	urlcolor=cyan		% color of external links
}
\renewcommand{\familydefault}{\sfdefault}
\usepackage{fancyhdr} % Custom headers and footers
\pagestyle{fancyplain} % Makes all pages in the document conform to the custom headers and footers
\fancyhead[L]{} % Empty left header
\fancyhead[C]{} % Empty center header
\fancyhead[R]{Daniel Pérez Ruiz} % My name
\fancyfoot[L]{} % Empty left footer
\fancyfoot[C]{} % Empty center footer
\fancyfoot[R]{\thepage} % Page numbering for right footer
%\renewcommand{\headrulewidth}{0pt} % Remove header underlines
\renewcommand{\footrulewidth}{0pt} % Remove footer underlines
\setlength{\headheight}{13.6pt} % Customize the height of the header

\usepackage{titlesec, blindtext, color}
\definecolor{gray75}{gray}{0.75}
\newcommand{\hsp}{\hspace{20pt}}
\titleformat{\chapter}[hang]{\Huge\bfseries}{\thechapter\hsp\textcolor{gray75}{|}\hsp}{0pt}{\Huge\bfseries}
\setcounter{secnumdepth}{4}
\usepackage[Lenny]{fncychap}


\begin{document}

	% Plantilla portada UGR
	\begin{titlepage}
\newlength{\centeroffset}
\setlength{\centeroffset}{-0.5\oddsidemargin}
\addtolength{\centeroffset}{0.5\evensidemargin}
\thispagestyle{empty}

\noindent\hspace*{\centeroffset}\begin{minipage}{\textwidth}

\centering
\includegraphics[width=0.9\textwidth]{logos/logo_ugr.jpg}\\[1.4cm]

\textsc{ \Large TRABAJO FIN DE GRADO\\[0.2cm]}
\small \textsc{DOBLE GRADO EN INGENIERÍA INFORMÁTICA Y MATEMÁTICAS}\\[1cm]

{\Huge\bfseries Simulador de Sistemas Concurrentes y Distribuidos \\}
\noindent\rule[-1ex]{\textwidth}{3pt}\\[3.5ex]
\end{minipage}

\vspace{1.5cm}
\noindent\hspace*{\centeroffset}
\begin{minipage}{\textwidth}
\centering

\textbf{Autor}\\ {Daniel Pérez Ruiz}\\[2.5ex]
\textbf{Director}\\ {Carlos Ureña Almagro}\\[1.2cm]

\begin{figure}[H]
  \centering
  \includegraphics[width=.3\textwidth]{logos/etsiit_logo.png}\hspace{0.2\textwidth}
  \includegraphics[width=.2\textwidth]{logos/ciencias-logo.png}
\end{figure}
%\includegraphics[width=0.3\textwidth]{logos/etsiit_logo.png}\\[0.1cm]
\textsc{Escuela Técnica Superior de Ingenierías Informática y de Telecomunicación}\\[0.1cm]
\textsc{Y}\\
\textsc{Facultad de Ciencias}\\
\textsc{---}\\
Granada, Junio de 2023
\end{minipage}
\end{titlepage}


	% Plantilla prefacio UGR
	\thispagestyle{empty}

\begin{center}
{\large\bfseries Simulador de Sistemas Concurrentes y Distribuidos. Lógica Temporal de Acciones }\\
\end{center}
\begin{center}
Daniel Pérez Ruiz\\
\end{center}

%\vspace{0.7cm}

\vspace{0.5cm}
\noindent\textbf{Palabras clave}: \textit{ágil, concurrencia, sistema, programa, lógica, instrucción, lenguaje, compilador, máquina virtual}
\vspace{0.7cm}

\noindent\textbf{Resumen}\\
Este trabajo se centra en el desarrollo de un compilador para un lenguaje de programación, nombrado Lamport en honor a Leslie Lamport y su influencia en el ámbito de los sistemas concurrentes y distribuidos, especialmente a través de su contribución con la Lógica Temporal de Acciones (TLA). Además este proyecto presenta un estudio detallado de TLA y un ejemplo práctico de su aplicación en la especificación de sistemas, proporcionando una perspectiva matemática y formal para la especificación y verificación de sistemas concurrentes en general. El objetivo principal, sin embargo, es la creación de una herramienta práctica que concretice los conceptos de concurrencia en una aplicación tangible y accesible para los usuarios.

El compilador de Lamport ha sido diseñado específicamente para simular sistemas concurrentes y distribuidos, permitiendo así a los usuarios interactuar directamente con las complejidades y retos inherentes a la programación en este campo. Desarrollado bajo una metodología ágil y utilizando las capacidades de los lenguajes C y C++, el compilador profundiza en la comprensión de aspectos clave como la sincronización de procesos o el manejo de bloqueos, a parte de ser un lenguaje totalmente funcional para programas secuenciales.

El trabajo también incluye una exploración de los principios fundamentales de los sistemas concurrentes y distribuidos, proporcionando el marco teórico esencial para comprender y utilizar eficazmente el compilador Lamport. Se examinan posibles ampliaciones y mejoras del compilador, incluyendo su eventual evolución hacia un modelo de Software as a Service (SaaS), la incorporación de mecanismos de sincronización más sofisticados y la expansión de la gramática del lenguaje.

\cleardoublepage

\begin{center}
	{\large\bfseries Concurrent and Distributed Systems Simulator. Temporal Logic of Actions}\\
\end{center}
\begin{center}
	Daniel Pérez Ruiz\\
\end{center}
\vspace{0.5cm}
\noindent\textbf{Keywords}: \textit{agile, concurrency, system, program, logic, instruction, language, compiler, virtual machine}
\vspace{0.7cm}

\noindent\textbf{Abstract}\\
This work focuses on the development of an compiler for a programming language, named Lamport in honor of Leslie Lamport and his influence in the field of concurrent and distributed systems, particularly through his contribution to Temporal Logic of Actions (TLA). Additionally, this project presents a detailed study of TLA and a practical example of its application in system specification, offering a mathematical and formal perspective for the specification and verification of concurrent systems in general. The main goal, however, is to create a practical tool that materializes the concepts of concurrency into a tangible and accessible application for users.

The Lamport compiler has been specifically designed to simulate concurrent and distributed systems, thus allowing users to directly interact with the complexities and challenges inherent in programming in this field. Developed under an agile methodology and utilizing the capabilities of C and C++ languages, the compiler deepens the understanding of key aspects such as process synchronization or deadlock handling, apart from being a fully functional language for sequential programs.

The work also includes an exploration of the fundamental principles of concurrent and distributed systems, providing the essential theoretical framework to effectively understand and use the Lamport compiler. Potential extensions and improvements of the compiler are examined, including its eventual evolution into a Software as a Service (SaaS) model, the incorporation of more sophisticated synchronization mechanisms, and the expansion of the language grammar.



\cleardoublepage

\thispagestyle{empty}

\noindent\rule[-1ex]{\textwidth}{2pt}\\[4.5ex]

D. \textbf{Tutora/e(s)}, Profesor(a) del ...

\vspace{0.5cm}

\textbf{Informo:}

\vspace{0.5cm}

Que el presente trabajo, titulado \textit{\textbf{Simulador de Sistemas Concurrentes y Distribuidos. Lógica Temporal de Acciones}},
ha sido realizado bajo mi supervisión por \textbf{Daniel Pérez Ruiz}, y autorizo la defensa de dicho trabajo ante el tribunal
que corresponda.

\vspace{0.5cm}

Y para que conste, expiden y firman el presente informe en Granada a Noviembre de 2023.

\vspace{1cm}

\textbf{El/la director(a)/es: }

\vspace{5cm}

\noindent \textbf{(nombre completo tutor/a/es)}

\chapter*{Agradecimientos}

Desde que era muy pequeño mi madre siempre me ha dicho que, si ves una mariposa blanca, eso es señal de suerte. Aunque no me caracterizo por ser una persona supersticiosa, siempre he convivido con esa inocencia cada vez que tenía ante mis ojos una de ellas, aleteando, al son de la libertad, iluminando en mí una sonrisa.

\vspace{0.2cm}

A lo largo de mi vida he tenido muchas ocasiones donde sentía que me habían abandonado. Y la realidad, es que siempre me han acompañado, aunque muchas veces no las viera. Ahí se encontraban, danzando al compás de todas las melodías y letras, de todas aquellas canciones que en tiempos pasados resonaron en mi cabeza.

\vspace{0.2cm} 

A mi madre, a mi hermana Irene, a mi padre. A mis amigos de toda una vida Adrián, Alberto, Elías, Juanmi, Juan, Rafa Aguilar y Rafa Barrales. A mis amigos de esta nueva vida Lucía, Martín, Mery, Jaime, Jose, Joseja y Pablo. No hay suficientes palabras en este mundo para describir lo inmensos que son. A mi amiga Candela, por haber por fin coincidido en este camino tras muchas historias y canciones, y que prometo recorrerlo bailando una vez más. A mi amiga Elena, por haber sido la mayor voz que me guió y me dio luz cuando no la había. A Coral, por haber venido casi al final con su deslumbrante forma de ser, brindándome un nuevo y prometedor comienzo. A mis dos grandes mentores: mi tutor, Carlos, por ser mi fuente de inspiración desde que coincidí con él; y JJ, la persona que me enseñó a cómo encontrarme si me perdía en el camino.

\vspace{0.2cm}
En este trabajo dejo atrás una muy buena parte de mi vida, donde cumplí el sueño de la infancia que me llevó hasta este mismo momento, el de descubrir y aprender todo lo \textit{immenso} que es el mundo de la informática, así como también lo es la vida misma. Esta parte de mí dejará de existir para siempre, pero ahora es mi turno de volar, al igual que todas las mariposas blancas que siempre estuvieron para mí.

\newpage
\vspace*{\fill}

\large
\noindent
$P := \textit{``Estás dispuesto a arriesgar.''}$
\newline
$Q := \textit{``Puedes crecer.''}$
\newline
$R := \textit{``Puedes dar lo mejor de ti.''}$
\newline
$S := \textit{``Puedes ser feliz.''}$
\newline
\newline

\noindent
\textit{Si} $\neg P \implies \neg Q$;
\newline
\textit{Si} $\neg Q \implies \neg R$;
\newline
\textit{Si} $\neg R \implies \neg S$;
\newline
\textit{Y si} $\neg S$, ¿qué te queda?
\newline



\vspace*{\fill}

	% Índice de contenidos
	\newpage
	\tableofcontents

	% Índice de imágenes y tablas
	\newpage
	\listoffigures

	% Si hay suficientes se incluirá dicho índice
	\listoftables 
	\newpage

	% Introducción 
	\chapter{\textbf{Introducción}}

\section{Motivación}
En el mundo actual, marcado por un avance tecnológico acelerado y una creciente digitalización, la importancia de los sistemas concurrentes y distribuidos se ha vuelto más evidente que nunca. Estos sistemas son fundamentales para el funcionamiento eficiente y óptimo de numerosas aplicaciones y servicios que forman la columna vertebral de la sociedad moderna. Desde el procesamiento de datos a gran escala hasta la computación en la nube, pasando por las redes de telecomunicaciones y la proliferación de la Internet de las Cosas (IoT), los sistemas concurrentes y distribuidos permiten el manejo y procesamiento simultáneo de enormes cantidades de información, asegurando la escalabilidad y la disponibilidad de los servicios.

En este entorno, la habilidad para diseñar, implementar y mantener sistemas que puedan manejar múltiples tareas de manera simultánea y distribuida es más que una habilidad técnica; es una necesidad vital. Con la integración creciente de la inteligencia artificial en diversos sectores, la eficiencia en la gestión de procesos concurrentes y la distribución efectiva de las cargas de trabajo no son solo deseables, sino esenciales para el avance y la innovación tecnológica. Además, la garantía de fiabilidad y seguridad en estos sistemas se convierte en un reto crítico, dado que cualquier fallo puede tener repercusiones significativas en un amplio espectro de actividades humanas.

En el grado de Ingeniería Informática de la Universidad de Granada, se dedica una asignatura completa al estudio de estos sistemas, denominada \textit{Sistemas Concurrentes y Distribuidos (SCD)}, donde se exploran tanto sus aspectos teóricos como prácticos. Esta asignatura no solo proporciona a los estudiantes una comprensión profunda de los fundamentos de los sistemas concurrentes y distribuidos, sino que también los introduce en el rigor matemático necesario para la verificación de sus propiedades utilizando un sistema lógico formal. Esta área de estudio es crucial, ya que asegura que los sistemas diseñados sean fiables, seguros y eficientes. Al enfrentarse a la complejidad inherente de estos sistemas, los alumnos aprenden a aplicar técnicas avanzadas de modelado y análisis.

La contrapartida es que su comprensión y aplicación práctica puede ser desafiante, especialmente para quienes se acercan por primera vez a la materia, y es por ello por lo que es crucial disponer de suficientes recursos y herramientas de apoyo que permitan alcanzar un aprendizaje íntegro y efectivo en este campo dinámico y en constante evolución.


\section{Objetivos del trabajo}
El propósito central de este proyecto es el desarrollo de un lenguaje de programación intuitivo y accesible , acompañado de un compilador específico, para facilitar el desarrollo y simulación de sistemas concurrentes. Este enfoque práctico busca proporcionar una alternativa a los lenguajes de programación de alto nivel convencionales, como C++, simplificando así el proceso de aprendizaje y enseñanza desde una perspectiva más accesible y transparente. Adicionalmente, este trabajo se adentra en un estudio sobre los fundamentos matemáticos subyacentes, centrándose en la Lógica Temporal de Acciones propuesta por Leslie Lamport. Este estudio es crucial para la especificación y verificación de los sistemas concurrentes, proporcionando una base teórica sólida que complementa y enriquece la comprensión práctica. Además, se explorarán los fundamentos detrás del procesamiento de lenguajes, esenciales para el diseño efectivo del lenguaje de programación propuesto. La combinación de estos enfoques prácticos y teóricos no solo busca lograr un entendimiento integral de los sistemas concurrentes y distribuidos, sino también fomentar una sinergia entre la teoría matemática y la aplicación práctica en el campo de la informática. A modo de resumen, los objetivos a cumplimentar son:

\begin{itemize}
    \item Desarrollar un marco teórico que abarque los conceptos clave de los sistemas concurrentes y distribuidos, proporcionando una comprensión clara y detallada de su funcionamiento, principios y desafíos. Incluir un estudio de los fundamentos matemáticos de las lógicas formales que sustentan las bases de la Lógica Temporal de Acciones, haciendo énfasis en su aplicación para la especificación y verificación de sistemas concurrentes. Se abordarán los principios fundamentales de la lógica y del tiempo, explorando cómo se integran y aplican en el contexto de la computación concurrente y distribuida.
    \item Llevar a cabo el desarrollo de un proyecto de software desde sus inicios utilizando metodologías de desarrollo ágil, poniendo el foco siempre en los clientes, que son las personas interesadas en el producto y quienes verdaderamente guían a un equipo hacia su progreso.
    \item Realizar un análisis y descripción detallada del lenguaje de programación a desarrollar, denominado \textit{Lamport}, tomando como punto de partida la sintaxis de los ejemplos de pseudocódigo de la asignatura \textit{Sistemas Concurrentes y Distribuidos} del grado de Ingeniería Informática de la Universidad de Granada.
    \item Implementar un compilador del código descrito, teniendo en cuenta las fases de análisis y optimización. Posteriormente, modelar una máquina virtual capaz de interpretar directamente el código generado, mostrando la ejecución resultante de las instrucciones al usuario, además de otorgarle transparencia en los sucesivos eventos que han ido ocurriendo por dentro, con el objetivo de garantizar un aprendizaje integral.
\end{itemize}

\section{Contenido del documento}\label{section:documentInfo}
Este documento está estructurado en 5 partes principales, cada una enfocada en aspectos distintos y fundamentales del proyecto. A continuación, se detalla el contenido de cada parte:

\begin{itemize}
    \item \textbf{Parte 1: Introducción y Motivación - El Origen y Propósito del Proyecto}. Esta sección explora la motivación detrás del proyecto, el problema que busca resolver y la metodología aplicada. Se detalla la planificación y evolución del proyecto, incluyendo un enlace al repositorio de GitHub para seguimiento práctico. Comprende los capítulos 1, 2 y 3.
    \item \textbf{Parte 2: Estado del Arte - El Contexto y la Relevancia Actual}. Se presenta un contexto histórico sobre la evolución de los sistemas concurrentes y distribuidos, destacando la contribución única del proyecto a esta área. Comprende el capítulo 4.
    \item \textbf{Parte 3: Fundamentos Teóricos y Matemáticos - La Base Conceptual del Proyecto}. Se realiza un análisis teórico completo de los sistemas concurrentes y distribuidos, abarcando desde la teoría matemática de sistemas lógicos formales, hasta los aspectos necesarios para la implementación práctica. Comprende los capítulos 5, 6 y 7.
    \item \textbf{Parte 4: Diseño e Implementación del Lenguaje de Programación Lamport - La Construcción de una Herramienta Innovadora}. Se tratan los aspectos técnicos del lenguaje Lamport, desde su diseño hasta la compilación, resaltando la importancia de los ejemplos prácticos para demostrar su eficacia. Comprende los capítulos 8, 9, 10, 11.
    \item \textbf{Parte 5: Conclusiones - Reflexiones y Perspectivas Futuras}. Se ofrecen reflexiones finales sobre los temas tratados y se sugieren direcciones futuras para la herramienta desarrollada, resumiendo los logros y aprendizajes clave del proyecto. Comprende el capítulo 12.

    
\end{itemize}

	% Descripción del problema y hasta donde se llega
	\chapter{\textbf{Descripción del problema}}\label{chapter:problema}

Este trabajo busca responder a las dos siguientes preguntas: ¿cómo se puede optimizar, a nivel teórico y práctico, el estudio e implementación de los conceptos de diseño y desarrollo de programas concurrentes? ¿Qué métodos existen para verificar formalmente las propiedades de dichos programas? Aunque existen innumerables maneras de abordar estos desafíos planteados, sólo el uso de una metodología clara, segura y correcta resulta verdaderamente adecuado.

\section{Metodología: Desarrollo Ágil}
La siguiente frase de Siegbert Tarrasch, quien fue uno de los mejores jugadores de ajedrez de todos los tiempos, es digna de mención: \textit{"La belleza de un movimiento no se refleja sólo en su apariencia, sino en el pensamiento detrás de él"}. No basta con tener acciones o hechos; es esencial que detrás de ellos exista una \textbf{idea}, un \textbf{problema} o simplemente una \textbf{pregunta} que sirva de base para construir paso a paso todos los objetivos que se deseen alcanzar. Del mismo modo que en el ajedrez cada movimiento es estratégico y sigue una lógica o un plan, en un proyecto que precise de la ingeniería informática cada paso dado debe estar fundamentado y orientado hacia la resolución del problema central.

En un proyecto de ingeniería informática, es esencial identificar claramente el problema que se desea resolver y la razón subyacente. Una vez definido, es importante garantizar que el proyecto en cuestión realmente aborde dicho problema \cite{jj-agile-objetivos}. La estrategia más efectiva consiste en dividir el problema en segmentos más manejables, a los que podemos referirnos como \textit{objetivos}. La recurrente mención de la palabra \textbf{"problema"} subraya su importancia central en la metodología, pues no se espera otra cosa que \textit{solucionar dicho problema}.

El desarrollo ágil surgió tras la redacción y firma del \textit{Manifiesto por el Desarrollo Ágil de Software} \cite{agile-manifest} por diecisiete expertos en programación. Con el término \textit{ágil} no se alude únicamente a una metodología para el desarrollo de proyectos que precisan de rapidez y flexibilidad, sino también una filosofía que implica una forma distinta de trabajar y de organizarse a la que predominaba anteriormente, denominada \textit{metodología en cascada}. Así pues, por \textit{ágil} se entiende una mentalidad que se aplica a todo el ciclo de vida del desarrollo de software, centrada en el cliente y en la mejora continua de productos mínimamente viables cada vez más complejos \cite{jj-agile-manifesto}.

Todo proyecto debe nacer de una motivación inicial, respondiendo a interrogantes como \textit{``por qué''}, \textit{``para qué''}, \textit{``para quién''} y \textit{``cómo''}. Si bien todas estas cuestiones son importantes, la penúltima destaca particularmente porque el éxito radica en satisfacer los \textit{deseos} y \textit{necesidades} de un grupo específico de clientes o usuarios unidos por una característica común: un \textit{problema}. 
Resolverlo implica practicar la \textbf{empatía} con ellos, esforzándose por comprender profundamente sus necesidades y determinar cómo satisfacerlas de manera óptima, aportando \textbf{valor} con los recursos disponibles. Las entrevistas personales o el seguimiento de las tendencias actuales pueden proporcionar la perspectiva adecuada.

De ahí surgen las \textbf{Historias de Usuario}, que proporcionan una explicación informal desde el punto de vista del usuario final y siempre situadas dentro del dominio del problema, de una funcionalidad del software que principalmente tendrá que ver con la lógica de negocio del proyecto.\cite{jj-design-thinking}. Una vez definidas, lo que queda es especificar los productos que se entregarán a los clientes, descritos a través de una secuencia de \textbf{hitos} o \textbf{milestones}. La esencia del desarrollo ágil radica en efectuar mejoras iterativas sobre el producto, contando siempre con la aprobación del usuario. En consecuencia, el avance del proyecto no es lineal, sino un proceso incremental.

En resumen, el desarrollo ágil supuso un gran cambio en el paradigma de la organización y planificación de proyectos de ingeniería informática, colocando al usuario y sus necesidades en el centro del proceso, guiando cada paso a través de la empatía y una comprensión profunda del problema. Con herramientas como las Historias de Usuario y la propia naturaleza de la metodología basada en un proceso iterativo e incremental, se busca proporcionar soluciones más adaptadas y flexibles de la ingeniería moderna, puesto que las necesidades de los clientes pueden ir variando con el tiempo. Es un enfoque que valora la colaboración y la resiliencia, con el objetivo de siempre aportar valor, abordando así los actuales y futuros desafíos del mundo tecnológico.

\section{Clientes}
Tras haber explorado la esencia del desarrollo ágil y su prioridad hacia el usuario y sus deseos, ahora se profundizará en el concepto de \textbf{cliente}. Los problemas y expectativas de los clientes o usuarios son lo que impulsan las decisiones y acciones del equipo de desarrollo. 

Para identificar y comprender los distintos usuarios del compilador de código, se empleó una metodología centrada en las personas. Se llevaron a cabo entrevistas individuales a dos grupos de personas \textit{reales} \footnote{Aunque no es estrictamente necesario utilizar personas reales para el análisis de un problema en el desarrollo de un proyecto, hacerlo añade una dimensión \textit{humana} que, en mi experiencia con otros proyectos, aumenta considerablemente las probabilidades de éxito.} vinculadas a la asignatura de \textit{Sistemas Concurrentes y Distribuidos}: estudiantes y profesorado. A ambos grupos se les formuló una serie de preguntas, que se presentan a continuación.

\noindent
Las preguntas planteadas al alumnado de la asignatura fueron:
\begin{itemize}
    \item \textit{¿Te resulta complicado entender los conceptos o fundamentos teóricos de la asignatura?}
    \item \textit{¿Eres capaz de resolver un problema que involucre concurrencia sin programarlo explícitamente?}
    \item \textit{¿Notas una diferencia significativa entre la teoría y práctica de la asignatura?}
    \item \textit{¿Crees que programar en el pseudocódigo específico de la asignatura te facilitaría su aprendizaje?}
\end{itemize}

\noindent
Las preguntas planteadas al profesorado fueron:
\begin{itemize}
    \item \textit{¿Consideras que la asignatura es difícil para los estudiantes?}
    \item \textit{¿Ves una diferencia significativa entre la teoría y la práctica de la asignatura?}
    \item \textit{¿Piensas que programar en el pseudocódigo específico de la asignatura ayudaría a los alumnos durante el curso? ¿Facilitaría tu labor docente al explicar los conceptos?}
\end{itemize}

Las dos últimas preguntas de cada grupo son similares, buscando comprobar la concordancia entre ambos puntos de vista. Las respuestas de los estudiantes indican que, aunque la asignatura posee una complejidad relativa en comparación con otras del grado, los conceptos en sí no son difíciles de asimilar. La verdadera barrera surge al aplicar estos conceptos en el diseño de sistemas concurrentes o al tratar de resolver problemas sin recurrir a un lenguaje de alto nivel, como C++. Los estudiantes generalmente consideran que se desenvuelven mejor en la práctica que en la teoría. Esta preferencia puede deberse a su familiaridad con el enfoque práctico adoptado en los primeros años del grado, con este lenguaje en particular. Por tanto, la respuesta a la última pregunta suele ser un \textit{sí} rotundo.

El profesorado, por su parte, no ve a su asignatura como particularmente complicada y no percibe una gran diferencia entre la teoría y práctica. Sin embargo, muestran empatía hacia los estudiantes, entendiendo que ellos puedan sentir una mayor complejidad. Coinciden en que disponer de un lenguaje de programación con la sintaxis propuesta en la asignatura podría facilitar una enseñanza más didáctica y accesible.

Con esta información no sólo obtenemos la motivación mencionada en la sección anterior, sino que también podemos identificar a los dos usuarios potenciales del compilador a desarrollar. A continuación, se describirá detalladamente el perfil de cada tipo de usuario.

\subsection{Tipo de usuario 1: Estudiante}
Se proporciona una descripción detallada de un perfil de usuario de tipo estudiante:

\begin{description}
    \item[Nombre:] Luis Martínez
    \item[Características Demográficas:] \hfill
        \begin{itemize}
            \item Edad: 19 años.
            \item Estudiante universitario de ingeniería informática, actualmente cursando la asignatura de \textit{Sistemas Concurrentes y Distribuidos}.
        \end{itemize}
    \item[Necesidades y Objetivos:] Luis aspira a aprobar la asignatura de \textit{Sistemas Concurrentes y Distribuidos} para avanzar en sus estudios. Además, busca comprender conceptos que sean importantes en futuras asignaturas relacionadas con su grado.
    
    \item[Habilidades Técnicas:] Posee habilidades de programación de nivel principiante a intermedio. Es probable que esta sea su primera experiencia con la implementación de programas no secuenciales.
    
    \item[Escenarios de Uso Comunes:] Luis utiliza el compilador para validar ejercicios específicos de la clase sobre sincronización de hebras o procesos.
    
    \item[Limitaciones:] Aunque Luis se siente cómodo programando en lenguajes de alto nivel como C++, enfrenta dificultades al comprender la sintaxis del pseudocódigo. Esto puede dificultar su capacidad para traducir rápidamente los conceptos teóricos en implementaciones prácticas utilizando pseudocódigo.
    
    \item[Expectativas:] Espera poder disponer de un lenguaje de pseudocódigo funcional, pudiendo añadir comentarios explicativos a cada línea de código para facilitar su comprensión. Además, desea que las ejecuciones de los programas se visualicen de forma clara, permitiéndole seguir el proceso paso a paso para consolidar su entendimiento de los conceptos teóricos.
\end{description}

\subsection{Tipo de usuario 2: Profesor}
Se proporciona una descripción detallada de un posible perfil de usuario de tipo profesor:

\begin{description}
    \item[Nombre:] Laura Ruiz
    \item[Características Demográficas:] \hfill
        \begin{itemize}
            \item Edad: 42 años.
            \item Profesora titular de la asignatura \textit{Sistemas Concurrentes y Distribuidos} en la facultad de Ingeniería Informática.
            \item Más de 10 años de experiencia docente en el campo de la informática.
        \end{itemize}
    \item[Necesidades y Objetivos:] Desea que sus estudiantes comprendan a fondo los conceptos y aplicaciones de los sistemas concurrentes y distribuidos. Busca herramientas y métodos que puedan hacer que la enseñanza sea más interactiva y efectiva.
    
    \item[Habilidades Técnicas:] Amplios conocimientos en programación, sistemas concurrentes, y pedagogía. Familiarizada con varios lenguajes de programación, incluido C++.
    
    \item[Escenarios de Uso Comunes:] Utilizar el compilador para demostrar ejemplos en clase, proponer ejercicios prácticos a los estudiantes y evaluar soluciones propuestas por ellos. Puede usarlo también para simular escenarios concurrentes y mostrar visualmente a los estudiantes cómo funcionan.
    
    \item[Limitaciones:] Prefiere que la herramienta tenga una interfaz amigable y sea intuitiva, ya que no desea invertir mucho tiempo en aprender a usarla. 
    
    \item[Expectativas:] Espera que el compilador permita explicaciones paso a paso y que pueda integrarse fácilmente con otros recursos didácticos. Le interesa que el pseudocódigo esté alineado con el contenido teórico de la asignatura, facilitando la transición entre teoría y práctica.
\end{description}

\newpage

\section{Historias de Usuario}

A través de las \textbf{Historias de Usuario}, se busca no solo definir qué es lo que el software debe hacer, sino también por qué es relevante hacerlo y cuál es el valor que se ofrece al usuario final. A continuación, se presentarán las Historias de Usuario identificadas para el desarrollo del compilador de código y cómo estas sientan las bases para los \textit{milestones} y Productos Mínimamente Viables del proyecto.

\begin{itemize}
    \item \textbf{Historia de usuario 1 (HU1):} Como estudiante de la asignatura, Luis quiere disponer de un compilador para escribir y ejecutar programas en el lenguaje de pseudocódigo propuesto en las transparencias, con el objetivo de practicar y verificar el funcionamiento de las soluciones a los ejercicios planteados.
    \item \textbf{Historia de usuario 2 (HU2):} Como profesora de la asignatura, Laura quiere contar con un compilador del lenguaje de las transparencias para poder explicar de forma más práctica los conceptos teóricos sin requerir un lenguaje de programación estándar como puede ser C++ o Java, además de poder revisar y analizar el código fuente desarrollado eventualmente por sus alumnos, y proporcionarles retroalimentación personalizada a los mismos.
\end{itemize}

\section{Milestones del proyecto: Productos Mínimamente Viables (PMV)}

Siguiendo la línea de las \textit{Historias de Usuario} y la importancia de establecer una comunicación efectiva con el usuario final, es el momento de definir productos entregables concretos que materialicen estas historias en el transcurso del desarrollo, que comúnmente se denomina \textbf{Productos Mínimamente Viables (PMV)}. Están diseñados no solo como representaciones tangibles del progreso, sino también como puntos de revisión donde se puede evaluar y adaptar el proyecto basándose en la retroalimentación aportada o por el propio cliente o por otro equipo de desarrollo. Los \textbf{milestones} o \textbf{hitos} sirven para secuenciar y organizar estos PMV en el ámbito del desarrollo ágil. A continuación, se enumeran los \textit{milestones} establecidos para este proyecto y cómo cada uno aporta al objetivo final del compilador de código.

\subsection{Milestone 0: Infraestructura del proyecto y definición del lenguaje.}\label{subsection:PMV0}
El primer hito del proyecto se centra en establecer una base sólida para su desarrollo. Comienza con una definición clara del problema, obtenida a través de entrevistas con usuarios potenciales, y la formulación detallada de este problema junto con la elaboración de las historias de usuario. Paralelamente, se lleva a cabo un estudio exhaustivo del lenguaje de pseudocódigo utilizado en la asignatura de Sistemas Concurrentes y Distribuidos, analizando sus componentes y definiendo su gramática utilizando la notación de Backus-Naur. En el aspecto técnico, se seleccionan cuidadosamente herramientas esenciales para el proyecto, incluyendo el lenguaje de programación para el compilador, un sistema de control de versiones, gestores de dependencias y tareas, herramientas de linting y ejecución de tests. Además, se configuran estas herramientas para garantizar la gestión eficiente de las dependencias, la calidad del código, y la automatización de tareas repetitivas como la compilación y la ejecución de tests.

\subsection{Milestone 1: Implementación del Analizador Léxico.}\label{subsection:PMV1}
Este hito se dedica a la construcción y validación del analizador léxico, una etapa crucial en el proceso de interpretación de cualquier lenguaje de programación. El foco principal es desarrollar una herramienta capaz de convertir una cadena de entrada en una serie de tokens identificables, que serán fundamentales para las etapas posteriores del compilador. Para lograrlo, se considera una herramienta generadora de analizadores léxicos que se ajuste a las especificaciones del lenguaje y las necesidades del proyecto. A continuación, se definen los tokens del lenguaje, incluyendo palabras reservadas, identificadores y operadores, junto con sus patrones de reconocimiento. La implementación del analizador léxico se realiza utilizando esta herramienta, asegurando su capacidad para reconocer y clasificar correctamente cada token. Además, se establece una tarea de generación del analizador léxico para facilitar su integración con el sistema, así como una tarea de limpieza de código objeto y ficheros compilados para mantener la organización y eficiencia del proyecto.

\subsection{Milestone 2: Implementación del Analizador Sintáctico.}\label{subsection:PMV2}
El segundo hito se enfoca en el diseño e implementación del analizador sintáctico, esencial para verificar la estructura correcta del código fuente en el lenguaje de pseudocódigo, según la gramática definida previamente. Este proceso incluye la definición de operadores, su aridad, asociatividad y orden de precedencia. Se selecciona una herramienta generadora de analizadores sintácticos y se ajusta la gramática del lenguaje. Tras implementar el analizador, se establece una estrategia para la recuperación de errores sintácticos y se define una tarea para su generación/compilación. Paralelamente, se desarrolla el Árbol Sintáctico Abstracto (AST), integrándolo con el analizador y definiendo una módulo para su gestión. Se implementan estructuras y funciones para la gestión de errores sintácticos y su integración con el analizador. Además, se llevan a cabo pruebas de integración y se desarrolla el programa principal del compilador, permitiendo la lectura de ficheros de texto plano, el procesamiento del pseudocódigo, la impresión del AST o los errores sintácticos detectados.

\subsection{Milestone 3: Implementación del Analizador Semántico.}\label{subsection:PMV3}
Este milestone se enfoca en el diseño e implementación del analizador semántico, una fase crítica que garantiza que el código fuente en el lenguaje de pseudocódigo no solo esté estructuralmente correcto (como asegura el analizador sintáctico) sino que también posea un significado lógico y coherente. El proceso incluye la descripción semántica del lenguaje y la implementación de una Tabla de Símbolos para gestionar identificadores y sus propiedades. Se desarrollan algoritmos clave para la resolución de nombres y la comprobación de tipos, asegurando que cada identificador se declare correctamente y que las operaciones entre diferentes tipos de variables sean válidas. Se establece una tarea de compilación para el módulo Analizador Semántico, facilitando su integración con otros componentes del compilador. Además, se realizan pruebas de integración para validar su funcionamiento en diversos escenarios de código. Se implementan estructuras y funciones para la gestión y reporte de errores semánticos, integrando el módulo de errores con el analizador para proporcionar retroalimentación específica al usuario. Finalmente, se habilita la función de análisis semántico dentro del flujo del compilador, proporcionando mensajes claros sobre el éxito del análisis o los errores detectados.

\subsection{Milestone 4: Análisis e implementación de la fase de generación de código intermedio.}\label{subsection:PMV4}
El cuarto hito se centra en el diseño e implementación de la fase de generación de código intermedio, un proceso clave que convierte el árbol de análisis sintáctico en una representación intermedia más cercana al código de máquina. Esta fase comienza con la selección de una representación intermedia adecuada para el lenguaje desarrollado y la definición de sus instrucciones, especificando operandos y usos. Se desarrolla un generador de código intermedio, con tablas para mapear variables, etiquetas y literales, y controladores para la generación y optimización del código. Paralelamente, se implementa un manejador de registros de eventos para documentar y ofrecer transparencia sobre las fases de análisis del compilador, incluyendo el análisis sintáctico, semántico y la generación del código intermedio. Se llevan a cabo pruebas de integración para asegurar la funcionalidad del módulo de generación de código intermedio y se habilita su función dentro del flujo general del compilador, junto con la generación de ficheros de registro de eventos para las distintas fases de análisis.


\subsection{Milestone 5: Análisis e implementación del compilador de código intermedio o máquina virtual}\label{subsection:PMV5}
El quinto hito del proyecto se centra en el diseño e implementación de un compilador para el código intermedio generado en las fases anteriores, que actúa como una máquina virtual. Esta máquina virtual está diseñada para ejecutar la representación intermedia, creando un entorno controlado que simula las operaciones de una máquina real de forma abstracta y agnóstica a la arquitectura del hardware. Para lograrlo, se implementan varios componentes clave: un esquema de traducción de direcciones virtuales a físicas, una abstracción de bloque de memoria, la memoria de la máquina virtual, un vector de registros de CPU, la CPU de la máquina virtual y un sistema operativo simulado para la máquina. Además, se aborda el tratamiento de excepciones como ZeroDivision e IndexOutOfBounds. La finalización de este milestone incluye realizar pruebas de integración exhaustivas para asegurar la funcionalidad y la integración adecuada de la máquina virtual con el resto del sistema.

\subsection{Milestone 6: Desarrollo teórico. Fundamentos matemáticos y tecnológicos}\label{subsection:PMV6}
El sexto milestone se enfoca en el desarrollo teórico, introduciendo los conceptos fundamentales de los sistemas concurrentes y distribuidos y realizando un estudio detallado de las lógicas formales, con énfasis en la Lógica Temporal de Acciones. Este enfoque teórico está diseñado para proporcionar una comprensión profunda de las bases matemáticas y lógicas necesarias para la verificación formal de sistemas concurrentes. Se explora la Lógica Temporal de Acciones, desarrollada por Leslie Lamport, para entender cómo se puede aplicar en el contexto de la verificación y el aseguramiento de la confiabilidad y coherencia en los sistemas concurrentes. El estudio se extiende a la explicación de cómo estos principios teóricos se integran y aplican en la práctica, resaltando su relevancia en el diseño, análisis y evaluación de sistemas que requieren una sincronización y coordinación efectivas. Este milestone no solo fortalece la base teórica del proyecto, sino que también proporciona a los usuarios las herramientas y conocimientos necesarios para abordar desafíos complejos en el campo de los sistemas concurrentes y distribuidos, preparándolos para enfrentar y resolver problemas prácticos con una sólida comprensión teórica.



	% Estado del arte
	% 	1. Crítica al estado del arte
	% 	2. Propuesta
	\chapter{\textbf{Estado del arte}}
La \textit{programación concurrente} sienta sus bases en la década de los 1960 con la aparición de los sistemas operativos multiprogramación, que pretendían resolver uno de los problemas más críticos que se daban en aquella época: el uso eficiente de los recursos hardware. Anteriormente, la CPU ejecutaba sólo un programa a la vez, quedando inactiva si se producían operaciones de entrada/salida subutilizando sus recursos, produciendo como resultado un rendimiento deficiente. Con la multiprogramación, varios programas podían residir simultáneamente en la memoria principal, permitiendo a la CPU ejecutar otro programa diferente mientras se gestionaba una interrupción de E/S. Esta gestión se facilitó gracias a la invención de los ``canales'', controladores de dispositivos que operaban de forma independiente. Esta revolucionaria técnica tuvo un gran impacto en algunos sistemas informáticos que se desarrollaban por aquel entonces, como es el caso de los ``mainframes'' \cite{TecnologiaInformaticaMainframe}.


La programación concurrente fue inicialmente motivo de preocupación de los diseñadores de sistemas operativos. Al final de esta década de 1960, los diseñadores de hardware desarrollaron máquinas multiprocesador, que supusieron todo un reto para la implementación de nuevos sistemas operativos adaptados a estos recursos, pero también definió una nueva oportunidad para que los desarrolladores de aplicaciones pudieran emerger y posicionarse en el mercado laboral como una profesión con grandes expectativas a futuro.



El primer gran reto de esta nueva técnica de programación fue resolver lo que se denomina como: \textit{el problema de la sección crítica} \cite{GeeksForGeeksCriticalSection}. Este desafío central se refiere al desarrollo de algoritmos efectivos para la sincronización de procesos concurrentes que requieren acceder a un recurso compartido. La importancia de resolver este problema radica en la necesidad de evitar conflictos y garantizar la coherencia de los datos cuando múltiples procesos intentan leer o modificar el mismo recurso simultáneamente. La correcta gestión de la sección crítica es crucial para asegurar que los sistemas concurrentes funcionen de manera fiable y eficiente. Para entender mejor y modelar este reto, se plantearon problemas teóricos como \textit{la cena de los filósofos} \cite{brosgol1996dining}, \textit{lectores y escritores} \cite{nithyasrikannathalreaders}, \textit{el barbero durmiente} \cite{SariSleepingBarber} o el \textit{problema de los fumadores} \cite{MTUSmokerProblem}.


Estos escenarios anteriormente mencionados han sido ampliamente estudiados, generando miles de artículos que proponen soluciones, mejoras, debates y nuevas implementaciones de primitivas de sincronización como los semáforos o los monitores, que simplificaban la tarea del programador. Paralelamente, los lenguajes de programación de alto nivel estaban proliferando y evolucionando. Un ejemplo notable es Simula \cite{Sklenar1997OOPSimula}, desarrollado en la década de 1960 por Ole-Johan Dahl y Kristen Nygaard. Considerado uno de los primeros lenguajes orientados a objetos, Simula también introdujo conceptos fundamentales que influirían en la simulación de concurrencia, aunque su principal aporte fue establecer las bases para la programación orientada a objetos, un paradigma que tendría un impacto significativo en el desarrollo futuro de lenguajes y técnicas de programación concurrente.


Al final de la década de 1970 y principios de 1980, el surgimiento de redes de ordenadores como ARPANET, que facilitó la computación en áreas amplias, y el desarrollo de tecnologías como Ethernet para redes locales, marcó el comienzo de una nueva era en la informática. Estos avances no solo transformaron la forma en que las computadoras se conectaban y comunicaban entre sí, sino que también ampliaron significativamente el campo de la programación concurrente. Mientras la programación concurrente se enfoca en gestionar múltiples procesos dentro de un mismo sistema, las redes de computadoras introdujeron el desafío de coordinar procesos que se ejecutan en diferentes máquinas físicas. Esto dio origen a lo que se conoce como \textit{programación distribuida}, un enfoque que extendía los principios de la programación concurrente a sistemas distribuidos cuya esencia reside en que los procesos interactúan entre ellos a través del envío de mensajes en vez de escribir y leer variables compartidas.


A medida que la programación concurrente y distribuida ganaban relevancia en la década de 1980 y 1990, su integración en los lenguajes de programación de alto nivel comenzó a materializarse de manera significativa. Un hito notable en este desarrollo fue el lenguaje de programación Ada \cite{burns1998concurrency}, introducido en 1983. Diseñado originalmente para satisfacer las exigentes necesidades de los proyectos de defensa y sistemas en tiempo real, Ada se destacó por su soporte nativo para la programación concurrente. Su enfoque en la seguridad, la fiabilidad y la gestión sofisticada de tareas concurrentes lo convirtió en una herramienta fundamental para aplicaciones críticas. Posteriormente, en la década de 1990, Java emergió como un lenguaje de programación influyente, llevando la programación concurrente a un público más amplio. Con su API de concurrencia, Java simplificó la gestión de hilos y procesos concurrentes, haciendo que la programación de este tipo fuera más accesible y manejable para los desarrolladores. Esta API proporcionó un conjunto robusto de herramientas y estructuras para el manejo eficiente de tareas concurrentes y sistemas distribuidos, reflejando la creciente demanda de aplicaciones que operaban en entornos de red y web. La evolución de la programación concurrente en Ada y Java no solo demuestra cómo se han adaptado los lenguajes de programación a desafíos técnicos complejos, sino también cómo han evolucionado para satisfacer las necesidades de aplicaciones modernas en diversos entornos de computación. Pero no solo los lenguajes en sí, sino también las herramientas para crearlos, evolucionaron significativamente durante este periodo. El desarrollo de herramientas avanzadas para generar compiladores facilitó la creación de lenguajes de programación más sofisticados, especialmente aquellos diseñados para la programación concurrente y distribuida. Estas herramientas permitieron a los diseñadores de lenguajes experimentar con nuevas construcciones y paradigmas, facilitando el desarrollo de lenguajes que pudieran manejar de manera más eficiente y segura las complejidades inherentes a la concurrencia y la distribución. 


Sin embargo, cada vez que se implementaban mejoras en la programación concurrente y distribuida, con nuevas ideas y hardware, la complejidad de los programas aumentaba exponencialmente. Esto llevó a un incremento similar en la dificultad de verificar su funcionamiento correcto. Poco después del origen de este modelo de programación, surgió la necesidad de introducir formalismos capaces de demostrar rigurosamente que un programa concurrente posee ciertas propiedades. En este contexto, las herramientas matemáticas formales, especialmente la lógica temporal, se convirtieron en indispensables. La lógica temporal, que se ocupa de los aspectos temporales del razonamiento en la lógica matemática, permitía especificar y razonar sobre el comportamiento de los programas a lo largo del tiempo, asegurando propiedades como la seguridad (nunca ocurrirá nada malo) y la vivacidad (realmente sucede algo bueno) en sistemas concurrentes.


Un desarrollo particularmente significativo en el campo de la verificación de sistemas fue la Lógica Temporal de Acciones (TLA) de Leslie Lamport, cuya vida y carrera han sido tan influyentes como sus logros académicos \footnote{Lamport fue también el desarrollador inicial de \LaTeX, el sistema de composición de documentos utilizado para redactar este informe. A lo largo de su carrera, ha sido galardonado con numerosos premios, incluyendo el prestigioso Premio Turing, por sus contribuciones a la informática teórica. Sus trabajos, incluyendo la TLA, han influenciado tanto la teoría como la práctica en el diseño y verificación de sistemas distribuidos y concurrentes.}. Nacido en 1941 en Nueva York, Lamport demostró un interés temprano en las matemáticas y la ciencia, lo que eventualmente lo llevó a obtener un Ph.D. en Matemáticas. Su transición hacia la computación fue motivada por un interés creciente en los problemas de sistemas distribuidos y la sincronización de procesos, un campo en el que se convertiría en un líder mundial.


Lamport, una figura destacada en la computación teórica, ha sido pionero en proponer metodologías rigurosas para el diseño y análisis de algoritmos en entornos complejos. La metodología de la TLA, que permite describir el comportamiento de un sistema en términos de estados y transiciones entre estos, facilita la verificación de propiedades críticas de sistemas complejos. Su visión era proporcionar una herramienta que no solo facilitara la especificación precisa de estos sistemas, sino que también permitiera su análisis y verificación de manera sistemática y confiable.


En la actualidad, la programación concurrente y distribuida se ha consolidado como un pilar fundamental en el panorama tecnológico. En el ámbito de la computación ubicua, donde múltiples dispositivos y sensores interactúan en tiempo real, es crucial mantener la precisión y la sincronización para ofrecer experiencias de usuario cohesivas y confiables. En el campo de la inteligencia artificial, la capacidad de manejar operaciones de procesamiento paralelo y distribuido resulta esencial para analizar grandes conjuntos de datos y ejecutar complejos algoritmos de aprendizaje automático. La corrección y eficiencia de los procesos concurrentes son vitales en áreas como el reconocimiento de patrones y el procesamiento de lenguaje natural. En el procesamiento multimedia, esta programación facilita la manipulación eficiente de contenidos de alta definición en tiempo real, siendo la precisión en la sincronización y la integridad de los datos fundamentales en aplicaciones que van desde la edición de vídeo hasta la transmisión en vivo. Estos ejemplos ilustran la omnipresencia y el impacto crucial de la programación concurrente y distribuida en la tecnología moderna, resaltando la necesidad de métodos rigurosos para su verificación fiable, lo que a su vez impulsa la innovación tecnológica.
	
	\chapter{Planificación}

\section{Metodología utilizada}


\section{Temporización}

\section{Seguimiento del desarrollo}


	% Análisis del problema
	% 1. Análisis de requisitos
	% 2. Análisis de las soluciones
	% 3. Solucion propuesta
	% 4. Análisis de seguridad
	\chapter{Bases teóricas}
 


	% Desarrollo bajo sprints: 
	% 	1. Permitir registros y login de usuarios
	% 	2. Desarrollo del sistema de incidencias
	% 	3. Desarrollo del sistema de denuncias administrativas y accidentes
	% 	4. Desarrollo del sistema de croquis
	%   5. Instalación de la aplicación de manera automática
	\chapter{Implementación}

La implementación del software se ha dividido en hitos. Estos, han sido definidos en Github
y cada uno de ellos contiene un grupo de \textit{issues} que se corresponden con las distintas
mejoras que se han ido incorporando al software a lo largo de su desarrollo.\\



	% Presupuesto

	% Conclusiones
	\chapter{Conclusiones y trabajos futuros}



	% Trabajos futuros


	
	\newpage
	\bibliography{bibliografia}
	\bibliographystyle{plain}
	
\end{document}

